\documentclass[letter,12pt]{article}
\usepackage[paperheight=27.94cm,paperwidth=21.59cm,bindingoffset=0in,left=3cm,right=2.0cm,top=3.5cm,bottom=2.5cm,headheight=200pt,headsep=1.0\baselineskip]{geometry}
\usepackage{graphicx,lastpage}
\usepackage{upgreek}
\usepackage{censor}
\usepackage[spanish,es-tabla]{babel}
\usepackage{pdfpages}
\usepackage{tabularx}
\usepackage{graphicx}
\usepackage{adjustbox}
\usepackage{xcolor}
\usepackage{colortbl}
\usepackage{rotating}
\usepackage{multirow}
\usepackage[utf8]{inputenc}
\usepackage{float}
\usepackage{hyperref}
\usepackage{listings}
\usepackage{xcolor}
\usepackage{booktabs}
\usepackage{amsmath}
\definecolor{codegreen}{rgb}{0,0.6,0}
\definecolor{codegray}{rgb}{0.5,0.5,0.5}
\definecolor{codepurple}{rgb}{0.58,0,0.82}
\definecolor{backcolour}{rgb}{0.95,0.95,0.92}
\lstdefinestyle{mystyle}{
    backgroundcolor=\color{backcolour},
    commentstyle=\color{codegreen},
    keywordstyle=\color{magenta},
    numberstyle=\tiny\color{codegray},
    stringstyle=\color{codepurple},
    basicstyle=\ttfamily\footnotesize,
    breakatwhitespace=false,
    breaklines=true,
    captionpos=b,
    keepspaces=true,
    numbers=left,
    numbersep=5pt,
    showspaces=false,
    showstringspaces=false,
    showtabs=false,
    tabsize=2
}
\lstset{style=mystyle}
\renewcommand{\tablename}{Tabla}
\usepackage{fancyhdr}
\pagestyle{fancy}
\begin{document}
\title{\Huge{Tarea 1 Sistemas Distribuidos} \\
  \large Plataforma de análisis de preguntas y respuestas en Internet}
\author{\\Dante Libiot, dante.libiot@mail.udp.cl\\ Benjamin Jimenez, benjamin.jimenez4@mail.udp.cl \\ \textbf{Profesor: } Nicolás Hidalgo}
\date{Abril, 2026}
\maketitle
\tableofcontents
\newpage
% ============================================================
\section{Introducción}
% ============================================================
El objetivo de este trabajo es diseñar e implementar un sistema de caché distribuido para
optimizar consultas geoespaciales sobre edificaciones en Santiago. El problema es simple:
cuando se consulta repetidamente la misma información, recalcular la respuesta cada vez
es ineficiente. La solución es guardar las respuestas en caché para que las consultas
repetidas se respondan desde memoria sin reprocesar los datos.

Se usa \textbf{Redis} como caché porque almacena datos en RAM, lo que lo hace
significativamente más rápido que releer un archivo o consultar una base de datos.
Además, al ser un servicio independiente, puede ser compartido por múltiples servidores,
lo que lo hace apropiado en un contexto distribuido.

% ============================================================
\section{Diseño de la Arquitectura}
% ============================================================
\subsection{Descripción General}

El sistema sigue una \textbf{arquitectura cliente-servidor de dos capas (2-tier)}.
El frontend actúa como \textit{thin client}: no tiene lógica de negocio, no sabe nada
de Redis ni del dataset, y su única función es mostrar información y mandar órdenes al
backend vía HTTP. El backend concentra toda la lógica: gestiona Redis, procesa consultas,
registra métricas y persiste el historial.

Esta separación tiene ventajas prácticas: el cliente es liviano (corre en cualquier
browser sin instalación), y si se modifica la lógica del caché o la política de evicción,
solo se toca el backend.

El flujo es el siguiente: el Generador de Tráfico manda consultas al Sistema de Caché.
Si la respuesta existe en Redis (\textit{hit}), se retorna de inmediato. Si no existe
(\textit{miss}), el Generador de Respuestas la calcula desde los datos en memoria, la
guarda en Redis y la retorna. Las métricas se registran en cada paso.

\subsection{Componentes Principales}

\subsubsection{Generador de Tráfico}
Simula consultas de empresas de logística sobre zonas de Santiago. Cada consulta combina:
\begin{itemize}
    \item \textbf{Zona geográfica} (Z1--Z5): seleccionada según la distribución configurada (uniforme o Zipf).
    \item \textbf{Tipo de consulta} (Q1--Q5): seleccionado de forma uniforme al azar.
    \item \textbf{Parámetros}: \texttt{confidence\_min} $\in [0,1]$ redondeado a 2 decimales para Q1--Q4; \texttt{bins} $\in \{5,10,20\}$ para Q5.
\end{itemize}

\subsubsection{Sistema de Caché (Redis)}
Recibe cada consulta, genera una clave única del formato \texttt{tipo:zona:parámetros},
y busca esa clave en Redis. Si existe, retorna la respuesta guardada. Si no, delega al
Generador de Respuestas y almacena el resultado con el TTL configurado.

\subsubsection{Generador de Respuestas}
Procesa las consultas que el caché no pudo responder, operando sobre los datos cargados
en memoria al inicio. No accede a disco durante la ejecución.

\subsubsection{Almacenamiento de Métricas}
Registra hits, misses, latencias, throughput y evictions. Al terminar cada experimento,
guarda un resumen completo en SQLite.

% ============================================================
\section{Implementación}
% ============================================================

\subsection{Dataset: Google Open Buildings}

Se usa el dataset Google Open Buildings, tile de Chile central, con $\sim$1.706.965
registros donde cada fila representa un edificio detectado por satélite. Los campos
relevantes son:

\begin{table}[H]
\centering
\begin{tabular}{lll}
\toprule
\textbf{Campo} & \textbf{Tipo} & \textbf{Descripción} \\
\midrule
\texttt{latitude}         & FLOAT & Latitud del centroide \\
\texttt{longitude}        & FLOAT & Longitud del centroide \\
\texttt{area\_in\_meters} & FLOAT & Área del edificio (m²) \\
\texttt{confidence}       & FLOAT & Certeza de detección {[}0, 1{]} \\
\bottomrule
\end{tabular}
\caption{Campos del dataset utilizados.}
\end{table}

\texttt{confidence} indica qué tan seguro está el modelo de que el objeto detectado es
un edificio. Actúa como filtro de calidad en Q1--Q4 y como variable de análisis en Q5.
Al iniciar el sistema se cargan todos los datos en memoria y se filtran los $\sim$29.170
edificios dentro de las cinco zonas, descartando el resto.

\subsection{Zonas Geográficas}

Cada zona es un rectángulo de coordenadas (\textit{bounding box}). Solo los edificios
dentro de esos límites se consideran para esa zona.

\begin{table}[H]
\centering
\begin{tabular}{llrrrr}
\toprule
\textbf{Zona} & \textbf{Nombre} & \textbf{lat\_min} & \textbf{lat\_max} & \textbf{lon\_min} & \textbf{lon\_max} \\
\midrule
Z1 & Providencia      & $-33{,}445$ & $-33{,}420$ & $-70{,}640$ & $-70{,}600$ \\
Z2 & Las Condes       & $-33{,}420$ & $-33{,}390$ & $-70{,}600$ & $-70{,}550$ \\
Z3 & Maipú            & $-33{,}530$ & $-33{,}490$ & $-70{,}790$ & $-70{,}740$ \\
Z4 & Santiago Centro  & $-33{,}460$ & $-33{,}430$ & $-70{,}670$ & $-70{,}630$ \\
Z5 & Pudahuel         & $-33{,}470$ & $-33{,}430$ & $-70{,}810$ & $-70{,}760$ \\
\bottomrule
\end{tabular}
\caption{Zonas geográficas con sus bounding boxes.}
\end{table}

\begin{table}[H]
\centering
\begin{tabular}{llr}
\toprule
\textbf{Zona} & \textbf{Nombre} & \textbf{Edificios} \\
\midrule
Z1 & Providencia      & 2.750  \\
Z2 & Las Condes       & 5.048  \\
Z3 & Maipú            & 13.250 \\
Z4 & Santiago Centro  & 4.350  \\
Z5 & Pudahuel         & 3.772  \\
\midrule
   & \textbf{Total}   & \textbf{29.170} \\
\bottomrule
\end{tabular}
\caption{Edificios por zona tras el filtrado inicial.}
\end{table}

\subsection{Tipos de Consultas (Q1--Q5)}

\subsubsection{Q1 --- Conteo de edificios}
Cuenta cuántos edificios de una zona tienen \texttt{confidence} $\geq$ umbral.\\
\textbf{Cache key:} \texttt{count:\{zona\}:conf=\{confidence\_min\}}

\subsubsection{Q2 --- Área promedio y total}
Calcula área promedio y total de los edificios filtrados por confianza.\\
\textbf{Cache key:} \texttt{area:\{zona\}:conf=\{confidence\_min\}}

\subsubsection{Q3 --- Densidad por km²}
Divide el conteo de Q1 por el área de la zona para comparar sectores de distinto tamaño.\\
\textbf{Cache key:} \texttt{density:\{zona\}:conf=\{confidence\_min\}}

\subsubsection{Q4 --- Comparación de densidad entre dos zonas}
Compara la densidad de dos zonas y retorna cuál tiene mayor concentración.\\
\textbf{Cache key:} \texttt{compare:density:\{zona\_a\}:\{zona\_b\}:conf=\{confidence\_min\}}

\subsubsection{Q5 --- Distribución de confianza}
Genera un histograma de los valores de \texttt{confidence} de una zona, agrupados en
\texttt{bins} intervalos. A diferencia de Q1--Q4, aquí \texttt{confidence} no es un
filtro sino el objeto de análisis.\\
\textbf{Cache key:} \texttt{confidence\_dist:\{zona\}:bins=\{bins\}}

% ============================================================
\section{Experimentos y Resultados}
% ============================================================

\subsection{Decisiones de Diseño}

\subsubsection{Arquitectura de Contenedores}

El sistema se despliega en \textbf{5 contenedores Docker} orquestados con un único
\texttt{docker-compose.yml} dentro de la misma red interna. Esto permite levantarlo
completo con un solo \texttt{docker-compose up} sin configuración manual.

\textbf{Redis} usa la imagen oficial Alpine. La política de evicción y el tamaño de
memoria son configurables en caliente desde el frontend, sin tocar el compose.

\textbf{Backend Servidor} (Python + FastAPI): núcleo del sistema. Procesa Q1--Q5,
maneja Redis, registra métricas y guarda el historial en SQLite via SQLAlchemy.

\textbf{Frontend Servidor} (HTML/JS + Bootstrap + Chart.js): permite configurar TTL,
política y memoria, ver métricas en vivo y comparar experimentos en gráficos. Todo
servido por FastAPI como archivos estáticos, sin framework JS.

\textbf{Backend Cliente} (Python + FastAPI): recibe la orden del frontend, avisa al
servidor qué distribución se usará y genera las consultas de forma asíncrona. Se optó
por tareas asíncronas dentro del mismo proceso en lugar de contenedores dinámicos,
ya que levantar un contenedor por consulta sería innecesariamente costoso.

\textbf{Frontend Cliente} (HTML/JS + Bootstrap): permite configurar distribución y
número de consultas, y muestra el progreso en tiempo real.

\textbf{Persistencia}: SQLite en volumen Docker. Cada experimento queda registrado con
todos sus parámetros y métricas, lo que permite construir gráficos comparativos de forma
acumulativa.

\subsubsection{Distribuciones de Tráfico}

Se implementan dos distribuciones para la selección de zonas, que representan dos
comportamientos opuestos:

\textbf{Distribución uniforme}: reparte el tráfico equitativamente entre las 5 zonas.
Las repeticiones de consultas ocurren por azar y el hit rate resultante es relativamente bajo.

\textbf{Distribución Zipf} (parámetro $s=1{,}5$): concentra el tráfico en pocas zonas,
replicando el comportamiento real donde ciertos sectores se consultan mucho más que otros.
Con $s=1{,}5$, Z1 recibe $\approx$57\% del tráfico y Z5 menos del 5\%. Esto genera
patrones de repetición claros que benefician al caché: las zonas populares permanecen
en Redis y el hit rate sube considerablemente respecto a la distribución uniforme.

\subsubsection{Políticas de Evicción}

Cuando el caché se llena, Redis elimina entradas según la política configurada. Se
evalúan dos:

\textbf{LFU (Least Frequently Used)}: elimina la entrada menos consultada en total.
Es la política esperada como mejor bajo Zipf, porque preserva exactamente las zonas
más populares y elimina las raramente consultadas.

\textbf{LRU (Least Recently Used)}: elimina la entrada que lleva más tiempo sin ser
consultada. Funciona bien bajo distribución uniforme donde no hay un patrón claro de
popularidad. Bajo Zipf puede eliminar entradas populares si no fueron consultadas
recientemente, aunque en el largo plazo son las más pedidas.

La hipótesis es: LFU $>$ LRU bajo distribución Zipf, y LRU $\approx$ LFU bajo
distribución uniforme.

\subsubsection{TTL}

El dataset es \textbf{estático}: el CSV nunca cambia durante la ejecución, por lo que
el TTL no sirve para mantener datos frescos. Su función real aquí es controlar cuánto
tiempo vive una entrada en el caché, lo que afecta directamente el hit rate.

Se intentó definir los valores de TTL como fracciones de la duración del experimento,
pero se descartó ese enfoque porque la duración varía según el TTL configurado: con
TTL alto hay más hits y cada hit involucra leer y deserializar $\sim$160 KB de Redis,
lo que agrega latencia real al experimento. Con TTL=0 todo se calcula desde pandas y
el experimento termina antes. Esta variación hace que el TTL como fracción de $T$ no
sea un parámetro controlable de forma reproducible.

Se optó por valores absolutos que permiten observar el efecto de distintas ventanas de
expiración de forma independiente de la duración del experimento:

\begin{table}[H]
\centering
\begin{tabular}{ll}
\toprule
\textbf{TTL (s)} & \textbf{Interpretación} \\
\midrule
0   & Sin caché: ninguna entrada se retiene, hit rate = 0\% \\
85  & Ventana corta: las entradas expiran frecuentemente \\
170 & Ventana media: las entradas duran varios minutos \\
200 & Ventana media-alta \\
600 & Ventana larga: las entradas sobreviven casi todo el experimento \\
\bottomrule
\end{tabular}
\caption{Valores de TTL evaluados (absolutos en segundos).}
\end{table}

TTL=0 actúa como caso borde inferior: no se guarda nada en Redis y el hit rate es 0\%.
Los valores intermedios permiten observar cómo aumenta el hit rate a medida que las
entradas viven más tiempo. Se observó empíricamente que con 18.000 consultas y TTL=600
el hit rate llegó al 85\% pero seguía creciendo al terminar el experimento, lo que
indica que el sistema no alcanzó el estado estacionario. Con un mayor número de consultas
se estima que el hit rate converge hacia $\sim$88\% bajo distribución Zipf con LFU.

\subsubsection{Espacio de Claves y Decisión de Precisión de \texttt{confidence\_min}}

Para que los tres tamaños de caché (50 MB, 200 MB, 500 MB) muestren diferencias
observables, es necesario cumplir dos condiciones: (1) que las consultas se repitan
lo suficiente como para generar hits, y (2) que el espacio total de claves supere
los 500 MB para que ningún tamaño pueda guardarlo todo.

La precisión de \texttt{confidence\_min} controla directamente cuántas claves distintas
existen, generando una tensión:

\begin{itemize}
    \item \textbf{Más decimales} $\rightarrow$ más claves posibles $\rightarrow$ las consultas casi no se repiten $\rightarrow$ hit rate $\approx$ 0\%.
    \item \textbf{Menos decimales} $\rightarrow$ pocas claves $\rightarrow$ el caché se llena rápido y todo es hit $\rightarrow$ las tres configuraciones de memoria son indistinguibles.
\end{itemize}

Con \textbf{4 decimales} hay más de 40.000 claves posibles; en 18.000 consultas casi
no hay repeticiones. Con \textbf{1 decimal} hay $\sim$550 claves; el caché se llena en
los primeros minutos y después da igual el tamaño. Con \textbf{2 decimales} se obtienen
101 valores posibles de \texttt{confidence\_min} y 4.055 claves únicas totales, punto
de equilibrio que combinado con Zipf genera hits relevantes sin saturar ningún tamaño.

\begin{table}[H]
\centering
\begin{tabular}{lrr}
\toprule
\textbf{Tipo de consulta} & \textbf{Claves posibles} & \textbf{Cálculo} \\
\midrule
Q1, Q2, Q3 & 1.515 & $5 \text{ zonas} \times 101 \times 3$ \\
Q4          & 2.525 & $25 \text{ pares} \times 101$ \\
Q5          & 15    & $5 \text{ zonas} \times 3 \text{ bins}$ \\
\midrule
\textbf{Total} & \textbf{4.055} & \\
\bottomrule
\end{tabular}
\caption{Universo de claves únicas con 2 decimales en \texttt{confidence\_min}.}
\end{table}

\subsubsection{Padding para Estresar la Memoria}

Con 4.055 claves y respuestas JSON de $\sim$100 bytes cada una, el espacio total sería
$\sim$400 KB, mucho menor que los 50 MB del caché más pequeño. Redis nunca evictaría
nada y los tres tamaños de memoria serían equivalentes.

Para forzar presión de memoria, cada respuesta almacenada en Redis incluye un campo
\texttt{\_pad} de 160 KB que \textbf{no se retorna al cliente}: se agrega al escribir
en Redis y se elimina al leer antes de responder. Esto simula que las respuestas
incluyen datos adicionales (geometrías, metadatos) sin afectar la latencia percibida.

Con esto, el espacio total posible es:
\[
4.055 \text{ claves} \times 160\ \text{KB} \approx \mathbf{619\ \text{MB}}
\]

Este valor supera los 500 MB, asegurando que ninguno de los tres tamaños pueda almacenar
todo el universo de respuestas:

\begin{table}[H]
\centering
\begin{tabular}{lrrr}
\toprule
\textbf{Tamaño} & \textbf{Entradas máx.} & \textbf{\% del espacio total} & \textbf{Evicciones esperadas} \\
\midrule
50 MB  & $\sim$312   & 7{,}7\%  & Muchas \\
200 MB & $\sim$1.250 & 30{,}8\% & Moderadas \\
500 MB & $\sim$3.125 & 77{,}1\% & Pocas \\
\bottomrule
\end{tabular}
\caption{Capacidad y presión esperada por configuración de caché.}
\end{table}

\subsubsection{Número de Consultas por Experimento}

Se eligieron \textbf{18.000 consultas} por dos razones. Primero, garantizar que el caché
más grande (500 MB, $\sim$3.125 entradas) se llene completamente y aún queden suficientes
consultas para que las políticas de evicción actúen: con 18.000 consultas quedan
$\sim$14.875 después de llenar el caché, una razón de $\approx 4{,}7:1$ entre fase
estable y fase de llenado. Segundo, a 10 ms por consulta esto equivale a $T = 180$ s
$\approx$ 3 minutos por experimento, tiempo razonable para ejecutar múltiples corridas.

\subsection{Métricas de Evaluación}

\begin{table}[H]
\centering
\begin{tabular}{ll}
\toprule
\textbf{Métrica} & \textbf{Definición} \\
\midrule
Hit rate         & $\dfrac{\text{hits}}{\text{hits} + \text{misses}}$ \\[1em]
Throughput       & Consultas por segundo \\[0.5em]
Latencia p50/p95 & Percentiles 50 y 95 del tiempo de respuesta \\[0.5em]
Eviction rate    & Evicciones por minuto \\[0.5em]
Cache efficiency & $\dfrac{\text{hits} \cdot t_{\text{cache}} - \text{misses} \cdot t_{\text{db}}}{\text{total}}$ \\
\bottomrule
\end{tabular}
\caption{Métricas de evaluación del sistema.}
\end{table}

Se usan percentiles (p50/p95) en lugar del promedio simple porque el promedio se distorsiona
con valores extremos, ocultando problemas reales de rendimiento.

\subsection{Resultados}

\subsubsection{Análisis de Variación de Tiempos de Vida (TTL)}

Para evaluar el impacto del tiempo de vida (TTL) en el rendimiento del sistema, se configuró
un escenario base utilizando la distribución de tráfico Zipf, una política de reemplazo LFU
y una memoria máxima de 500MB. Se enviaron 18.000 consultas, variando el TTL y monitoreando
el comportamiento asintótico de las métricas.

\begin{itemize}
    \item \textbf{Hit rate y eficiencia del caché}\\
    El Hit Rate exhibe el comportamiento logarítmico esperado frente al incremento del TTL,
    dictado por la localidad temporal dentro de la distribución Zipf. Como se observa en la
    Figura~\ref{fig:ttl_hitrate}, partiendo de un 0\% para $TTL=0s$, el sistema escala
    progresivamente a un 73\% con $TTL=85s$, alcanzando un 85\% con $TTL=600s$.

    La curva presenta rendimientos marginales decrecientes en la eficiencia del caché. Además,
    al finalizar la inyección de las 18.000 consultas bajo $TTL=600s$, el \textit{Hit Rate}
    continuaba aumentando, evidenciando que el sistema se encontraba aún en su fase de
    calentamiento. Se proyecta que, con un volumen mayor de peticiones, la métrica convergería
    asintóticamente hacia un $\sim$88\%. Esta tendencia se corrobora con la eficiencia del caché
    (Figura~\ref{fig:ttl_eficiencia}): el sistema se vuelve eficiente con TTL alrededor de
    $170$--$200$ segundos, y desde ese punto el beneficio adicional al aumentar el TTL es
    casi nulo.

    \begin{figure}[H]
    \centering
    \begin{minipage}{0.48\textwidth}
        \centering
        \includegraphics[width=\linewidth]{ZIPF_HitRate-ttl.jpeg}
        \caption{Evolución del Hit Rate según el TTL.}
        \label{fig:ttl_hitrate}
    \end{minipage}\hfill
    \begin{minipage}{0.48\textwidth}
        \centering
        \includegraphics[width=\linewidth]{ZIPF_Cache_eficiencia-ttl.jpeg}
        \caption{Eficiencia de la caché variando el TTL.}
        \label{fig:ttl_eficiencia}
    \end{minipage}
    \end{figure}

    \item \textbf{Latencia (p50 y p95)}\\
    La latencia mediana ($p50$), presente en la Figura~\ref{fig:ttl_p50}, decrece a medida
    que aumenta el TTL, pasando de 0{,}66 ms ($TTL=0s$) a 0{,}34 ms ($TTL=600s$).

    Sin embargo, la latencia del percentil 95 ($p95$) presenta un fenómeno particular
    (Figura~\ref{fig:ttl_p95}). Con $TTL=85s$, el $p95$ alcanza los 1{,}53 ms, superando
    marginalmente al sistema sin caché ($TTL=0s$, 1{,}50 ms). Esta anomalía se explica por la
    expiración temprana de claves: las entradas se caducan constantemente durante la ejecución,
    generando \textit{cache misses} sobre datos previamente cacheados. Al extender el TTL,
    las expiraciones disminuyen drásticamente, estabilizando el $p95$ en 1{,}26 ms ($TTL=600s$).

    \begin{figure}[H]
    \centering
    \begin{minipage}{0.48\textwidth}
        \centering
        \includegraphics[width=\linewidth]{ZIPF_Latencia_P50-ttl.jpeg}
        \caption{Latencia mediana (p50) vs TTL.}
        \label{fig:ttl_p50}
    \end{minipage}\hfill
    \begin{minipage}{0.48\textwidth}
        \centering
        \includegraphics[width=\linewidth]{ZIPF_Latencia_P95-ttl.jpeg}
        \caption{Latencia de cola (p95) vs TTL.}
        \label{fig:ttl_p95}
    \end{minipage}
    \end{figure}

    \item \textbf{Throughput}\\
    El Throughput (Figura~\ref{fig:ttl_throughput}) experimenta un incremento moderado y
    progresivo (de 74 a 79 req/s) a medida que aumenta el TTL. Esta diferencia de $\sim$5\%
    indica que el límite superior está acotado por el cliente generador de tráfico, en
    específico por el \textbf{sleep} de 10 ms entre peticiones, lo que convierte a esta
    métrica en un factor del lado del cliente más que en un límite del servidor.

    \begin{figure}[H]
    \centering
    \includegraphics[width=0.6\textwidth]{ZIPF_Throughput-ttl.jpeg}
    \caption{Throughput del sistema frente a variaciones de TTL.}
    \label{fig:ttl_throughput}
    \end{figure}
\end{itemize}

\subsubsection{Comparativa de Política de Reemplazo}

En esta fase, se sometió al sistema a un escenario con diferentes políticas de reemplazo
(LFU y LRU) bajo ambas distribuciones de tráfico y una memoria de 500MB. El objetivo es
evaluar cómo cada política maneja la retención de las claves más populares cuando el espacio
en caché se agota.

\begin{itemize}
    \item \textbf{Hit rate y eficiencia del caché}\\
    Como se observa en la Figura~\ref{fig:pol_hitrate}, la política \textbf{LFU} se posiciona
    como la más óptima, alcanzando un acierto cercano al 88\%, mientras que \textbf{LRU}
    alcanza un $\sim$85\%. LFU es superior porque basa su decisión en la frecuencia histórica,
    protegiendo de manera constante las consultas a las zonas más densas. Por otro lado, LRU
    es marginalmente más susceptible a la polución de caché provocada por ráfagas de consultas
    atípicas en momentos específicos.

    La eficiencia del caché (Figura~\ref{fig:pol_eficiencia}) presenta lo inverso: \textbf{LRU}
    supera a \textbf{LFU}. Aunque LFU maximiza los aciertos históricos, carece de un mecanismo
    natural de envejecimiento temporal. Esto provoca que claves con ráfagas de popularidad en
    etapas tempranas permanezcan en memoria indefinidamente, incluso si dejaron de ser relevantes.
    LRU permite limpiar pasivamente la memoria de datos inactivos, lo que le otorga una rotación
    de claves más saludable y un uso más eficiente del espacio frente a recursos limitados.

    \begin{figure}[H]
    \centering
    \begin{minipage}{0.48\textwidth}
        \centering
        \includegraphics[width=\linewidth]{HitRate-Politica.jpeg}
        \caption{Comparativa de Hit Rate según política de evicción.}
        \label{fig:pol_hitrate}
    \end{minipage}\hfill
    \begin{minipage}{0.48\textwidth}
        \centering
        \includegraphics[width=\linewidth]{eficiencia_cache-politica.jpeg}
        \caption{Eficiencia de caché por política de reemplazo.}
        \label{fig:pol_eficiencia}
    \end{minipage}
    \end{figure}

    \item \textbf{Latencia (p50 y p95)}\\
    En cuanto a la latencia mediana ($p50$), Figura~\ref{fig:pol_p50}, ambas políticas
    operan por encima de los 0{,}7 ms; sin embargo, la latencia de \textbf{LRU} es
    consistentemente mayor que la de \textbf{LFU}. Esta diferencia es consecuencia directa
    del mayor hit rate de LFU: al retener de manera más robusta las \textit{hot-keys}
    históricas, resuelve una mayor proporción de consultas desde Redis sin delegar al backend.

    Este comportamiento se replica en $p95$ (Figura~\ref{fig:pol_p95}). Dado que LRU es más
    susceptible a desalojar claves temporalmente inactivas pero de alta recurrencia global,
    LFU logra contener mejor los picos de latencia de cola.

    \begin{figure}[H]
    \centering
    \begin{minipage}{0.48\textwidth}
        \centering
        \includegraphics[width=\linewidth]{Latencia_P50-politica.jpeg}
        \caption{Política en la latencia mediana (p50).}
        \label{fig:pol_p50}
    \end{minipage}\hfill
    \begin{minipage}{0.48\textwidth}
        \centering
        \includegraphics[width=\linewidth]{Latencia_P95-politica.jpeg}
        \caption{Política en la latencia de cola (p95).}
        \label{fig:pol_p95}
    \end{minipage}
    \end{figure}

    \item \textbf{Throughput}\\
    Las políticas LFU y LRU logran estabilizarse en torno a las 75 req/s
    (Figura~\ref{fig:pol_throughput}), maximizando la capacidad del sistema bajo las
    limitantes del cliente. La tasa de aciertos se relaciona directamente con la
    capacidad total de concurrencia.

    \begin{figure}[H]
    \centering
    \includegraphics[width=0.6\textwidth]{throughput-politica.jpeg}
    \caption{Throughput según heurística.}
    \label{fig:pol_throughput}
    \end{figure}
\end{itemize}

\subsubsection{Análisis de Variación de Tamaño de Memoria}

Para evaluar el impacto de la memoria en el rendimiento, se configuró un escenario base con
$TTL=600s$, promediando LFU y LRU y ambas distribuciones. Se enviaron 18.000 consultas
variando el límite máximo de memoria en 50MB, 200MB y 500MB.

\begin{itemize}
    \item \textbf{Hit rate y eficiencia del caché}\\
    Los resultados demuestran una relación directa entre la capacidad de memoria y la tasa de
    aciertos (Figura~\ref{fig:mem_hitrate}). La configuración más restrictiva (50MB) presenta
    el \textit{Hit Rate} más bajo porque el espacio es insuficiente para albergar el conjunto
    de trabajo principal, forzando reemplazos constantes. Al incrementar la memoria a 200MB
    y 500MB, el \textit{Hit Rate} mejora progresivamente, ya que el sistema retiene las
    consultas populares sin recurrir a reemplazos agresivos.

    Esta dinámica se refleja en la eficiencia del caché (Figura~\ref{fig:mem_eficiencia}).
    La configuración de 50MB muestra un rendimiento inferior debido a la alta rotación de
    datos. Las capacidades mayores logran una eficiencia superior al estabilizar el
    almacenamiento de las claves más consultadas.

    \begin{figure}[H]
    \centering
    \begin{minipage}{0.48\textwidth}
        \centering
        \includegraphics[width=\linewidth]{Hit_memoria.jpeg}
        \caption{Hit Rate según la memoria.}
        \label{fig:mem_hitrate}
    \end{minipage}\hfill
    \begin{minipage}{0.48\textwidth}
        \centering
        \includegraphics[width=\linewidth]{eficiencia_cache_memoria.jpeg}
        \caption{Eficiencia de la caché variando el tamaño de memoria.}
        \label{fig:mem_eficiencia}
    \end{minipage}
    \end{figure}

    \item \textbf{Latencia (p50 y p95)}\\
    La latencia $p50$ (Figura~\ref{fig:mem_p50}) es más alta con 50MB. Al ampliar la memoria
    a 200MB la latencia mediana cae drásticamente, y de 200MB a 500MB se estabiliza, dado que
    la mayoría de las consultas ya son hits en ambos casos.

    El mismo patrón se observa en $p95$ (Figura~\ref{fig:mem_p95}): la restricción de 50MB
    genera los picos de latencia más altos del experimento. Los escenarios de 200MB y 500MB
    absorben el tráfico de forma más estable al reducir la frecuencia de desalojos.

    \begin{figure}[H]
    \centering
    \begin{minipage}{0.48\textwidth}
        \centering
        \includegraphics[width=\linewidth]{latencia_p50_memoria.jpeg}
        \caption{Latencia mediana (p50) según la memoria.}
        \label{fig:mem_p50}
    \end{minipage}\hfill
    \begin{minipage}{0.48\textwidth}
        \centering
        \includegraphics[width=\linewidth]{latencia_p95_memoria.jpeg}
        \caption{Latencia de cola (p95) según la memoria.}
        \label{fig:mem_p95}
    \end{minipage}
    \end{figure}

    \item \textbf{Throughput}\\
    La configuración de 50MB presenta el throughput más bajo al enfrentar una mayor tasa de
    fallos de caché (Figura~\ref{fig:mem_throughput}). Aumentar la capacidad a 200MB y 500MB
    incrementa el volumen de consultas procesadas por segundo.

    \begin{figure}[H]
    \centering
    \includegraphics[width=0.6\textwidth]{throughput_memoria.jpeg}
    \caption{Throughput según memoria.}
    \label{fig:mem_throughput}
    \end{figure}
\end{itemize}

% ============================================================
\section{Discusión}
% ============================================================

\begin{itemize}
    \item \textbf{¿Cómo varía la tasa de aciertos (hit rate) y de fallos (miss rate) entre una distribución y otra?}\\
    La variación entre ambas distribuciones es drástica y está dictada por la probabilidad de
    recurrencia de las consultas. Bajo la \textbf{distribución Zipf}, el tráfico presenta alta
    localidad temporal, donde un grupo pequeño de claves concentra la mayoría de las peticiones.
    Esto permite que el sistema alcance tasas de aciertos superiores al 85\%, minimizando el
    \textit{Miss Rate}. Por el contrario, bajo una \textbf{distribución Uniforme}, las peticiones
    se dispersan equitativamente a lo largo de todo el espacio de claves posibles, provocando que
    a una consulta se le expire su TTL antes de repetirse o que sea desalojada por falta de
    memoria, lo que hace que el \textit{Hit Rate} decaiga y el \textit{Miss Rate} incremente.

    \item \textbf{¿Qué distribución de tráfico genera un mayor estrés o beneficio para la caché y por qué?}\\
    La distribución Uniforme genera mayor estrés. Al provocar una alta tasa de fallos constantes,
    la caché entra en un ciclo donde se ve obligada a escribir, desalojar y reescribir datos ya
    cacheados, transfiriendo toda la carga de procesamiento al backend y saturando la CPU.

    En contraparte, la distribución Zipf genera un beneficio para la caché. Al concentrar el
    tráfico en consultas populares, se cumple el propósito de diseño: absorber la carga pesada
    resolviendo las peticiones en memoria, protegiendo al servidor de cálculos redundantes y
    maximizando el throughput.

    \item \textbf{¿Qué implicaciones tienen estos resultados para un escenario de aplicación real?}\\
    Los resultados demuestran que el valor de un sistema de caché depende del comportamiento del
    problema. En escenarios reales de análisis de densidad en Santiago, el tráfico es naturalmente
    asimétrico (Zipfian): las zonas comerciales o centros urbanos densamente poblados reciben
    muchas más consultas que las zonas periféricas.

    Los datos comprueban que, dado que el mundo real no se comporta de manera uniforme, un diseño
    con caché limitada (e.g., 200MB) pero con políticas de reemplazo adecuadas (LRU o LFU) es
    capaz de absorber más del 85\% del tráfico, reduciendo los costos operativos y garantizando
    tiempos de respuesta óptimos.
\end{itemize}

% ============================================================
\section{Conclusiones}
% ============================================================

El presente informe demuestra que la integración de un caché distribuido en memoria (Redis) es
una decisión que optimiza sistemas logísticos que procesan grandes volúmenes de datos geoespaciales.
Al simular patrones de tráfico asimétricos realistas (distribución Zipf), se comprobó que el
sistema es capaz de absorber de manera eficiente la carga y mitigar el estrés computacional
del backend.

A partir del análisis de las métricas de rendimiento, se llegaron a las siguientes conclusiones:

\begin{itemize}
    \item \textbf{Umbral de rentabilidad del TTL:} Un tiempo de vida cercano a los \textbf{200 segundos}
    representa un punto óptimo. Superar este umbral aporta beneficios insignificantes, mientras que
    un valor inferior induce una alta tasa de expiraciones tempranas, penalizando la latencia.

    \item \textbf{Política de Reemplazo:} Aunque la política LFU maximiza el acierto histórico,
    la \textbf{política LRU} es también una opción recomendada. Su mecanismo de limpieza pasiva
    previene la retención indefinida de claves con popularidad pasajera, asegurando un uso eficiente
    de la memoria frente a variaciones temporales en el patrón de acceso.

    \item \textbf{Dimensionamiento:} No es necesario almacenar la totalidad del dataset en caché.
    Una \textbf{memoria de 200MB} demostró ser suficiente para albergar las consultas frecuentes,
    estabilizando el sistema sin incurrir en costos innecesarios de infraestructura.
\end{itemize}

En síntesis, una arquitectura de caché dimensionada al \textit{working set} real permite construir
una plataforma altamente escalable, capaz de soportar alta concurrencia y mantener tiempos de
respuesta bajos.

% ============================================================
\section{Referencias}
% ============================================================

\begin{itemize}
    \item Google Open Buildings Dataset: \url{https://sites.research.google/gr/open-buildings/}
    \item Redis Documentation: \url{https://redis.io/docs/}
    \item Docker Documentation: \url{https://docs.docker.com/}
\end{itemize}

\section*{Links}
\begin{enumerate}
    \item Repositorio de Github:
    \item Video de funcionamiento en Youtube:
\end{enumerate}

\end{document}
